\documentclass[11pt]{article}
\usepackage[margin=1in]{geometry}
\usepackage{amsmath,amssymb,amsthm}
\usepackage[T1]{fontenc}
\usepackage[hidelinks]{hyperref}
\newtheorem{proposition}{Proposition}
\newtheorem{corollary}[proposition]{Corollary}
\newcommand{\E}{\mathbb E}
\newcommand{\C}{\mathbb C}
\newcommand{\R}{\mathbb R}
\title{Exact outward transfers: Gaussian, Gamma and simplex moments}
\author{Independent internal research note}
\date{16 November 2026}
\begin{document}
\maketitle
\noindent\textbf{Scope.}
This is an elementary transfer argument and an application assessment.
No Factorial, Image, Dixmier, center-focus or meridional-rank conjecture is
claimed solved. No novelty claim or registry promotion is made. The input
Gaussian decision theorem has the assurance limits of its canonical proof.
No agents were messaged for this assessment.

\section{A complete Gamma--Dirichlet--Gaussian transfer}

Fix positive integers $k_1,\ldots,k_r$ and put
$\alpha_i=k_i/2$ and $A=\sum_i\alpha_i$.
Let $X_{ij}$ be independent standard real Gaussians and define
\[
 Y_i=\frac12\sum_{j=1}^{k_i}X_{ij}^2,\qquad
 R=\sum_iY_i,\qquad T_i=Y_i/R.
\]
The $Y_i$ are independent $\operatorname{Gamma}(\alpha_i,1)$ variables.
Moreover $R$ and $T$ are independent, with distributions
$\operatorname{Gamma}(A,1)$ and $\operatorname{Dir}(\alpha_1,\ldots,\alpha_r)$.
These distributional facts are classical; for example see
\href{https://www2.stat.duke.edu/~st118/sta941/dirichlet.pdf}{the Duke Dirichlet notes}.

Let $p$ be a polynomial on $\Delta=\{t_i\ge0,\sum_i t_i=1\}$, and
choose $D\ge\deg p$. Its homogeneous lift is
\[
 F(Y)=R^D p(Y/R).
\]
For a marker $q$ of degree at most $e$, similarly put
$G(Y)=R^e q(Y/R)$.
Both expressions are polynomials. They depend only on the polynomial
functions on the simplex when $D,e$ are fixed, because
$\sum_i(Y_i/R)=1$. With $(a)_n=a(a+1)\cdots(a+n-1)$,
\[
 \E[G(Y)F(Y)^m]=(A)_{e+Dm}\E[q(T)p(T)^m].                    \tag{1}
\]
Every factor $(A)_n$ is a nonzero rational number.

\begin{proposition}
For the indicated lift $F$, the following properties transfer in both
directions between $p$ under the Dirichlet functional and $F$ under the
standard Gaussian functional on all $X_{ij}$:
\begin{enumerate}
\item vanishing of every positive pure moment;
\item vanishing of all sufficiently large pure moments;
\item eventual marked vanishing for every fixed polynomial marker.
\end{enumerate}
The equivalence retains \emph{all} Gaussian markers, not only radial ones.
\end{proposition}
\begin{proof}
Equation (1) proves the pure assertions and the implication from Gaussian
safety to Dirichlet safety. Conversely, average any Gaussian polynomial
marker over $K=\prod_iO(k_i)$. Its expectation against $F^m$ is unchanged,
and the average is a polynomial in $Y_1,\ldots,Y_r$. Decompose it into
finitely many homogeneous pieces $H_e(Y)$, and apply (1) with
$q_e(T)=H_e(T)$ to each piece. Eventual Dirichlet marked vanishing gives
eventual vanishing of their finite sum. The orthogonal invariants here
are elementary radial invariants in each separate vector block; there
are no cross-block scalar products because the group factors act
independently.
\end{proof}

The same averaging argument, without homogenization, transfers all-marker
safety between any polynomial $f(Y)$ under independent Gamma variables
and $f(Y(X))$ under Gaussians. All arguments are coefficientwise and use
rational moment factors, so they extend to commutative $\mathbb Q$-algebras
for the corresponding rationally defined moment functionals. This does
not confer an executable input representation on an arbitrary coefficient ring.

There is an explicit bad-marker translation. A Gaussian monomial containing
an odd exponent in any coordinate has zero marked moments against $F^m$,
so a monomial bad marker must have the form
$Q(X)=\prod_{i,j}X_{ij}^{2b_{ij}}$. Put
$b_i=\sum_j b_{ij}$ and $b=\sum_i b_i$. Orthogonal averaging gives
\[
 \operatorname{Av}_K Q=\kappa Y_1^{b_1}\cdots Y_r^{b_r},\qquad
 \kappa=\prod_i
 \frac{2^{b_i}\prod_j(1/2)_{b_{ij}}}{(\alpha_i)_{b_i}}\ne0.
\]
Consequently
\[
 \E[Q(X)F(Y(X))^m]
   =\kappa(A)_{b+Dm}\E[T_1^{b_1}\cdots T_r^{b_r}p(T)^m].       \tag{2}
\]
Thus a supplied Gaussian monomial bad-marker certificate translates to
an original Dirichlet monomial bad marker. There is no inference that a
finite nonzero prefix alone proves infinite survival.

\begin{corollary}
The repository's general pointwise Gaussian safety producer over
$\mathbb Q$ or a simple number field transfers to polynomial phases under
every half-integer Dirichlet law, including uniform measure on a simplex
($k_i=2$ for every $i$), and to independent half-integer Gamma laws.
\end{corollary}
The corollary uses the existing producer and (1)--(2); it is not a new
implementation or a minimal-cutoff claim. A Gaussian bound
$m>C\deg_X Q+E$ yields the sufficient Dirichlet bound $m>2Ce+E$ by
homogenizing a degree-$e$ marker. These are exact algebraic moment
identities, not convergence assertions about $\E e^{tF}$ for nonzero $t$.

\section{The Factorial Conjecture and polynomial automorphisms}

For $k_i=2$, the $Y_i$ are independent exponentials and
\[
       \E[Y_1^{a_1}\cdots Y_r^{a_r}]=a_1!\cdots a_r!.
\]
This is exactly the factorial functional $\mathcal L$. Hence the Gaussian
machinery applies to specified factorial-moment instances after a polynomial
substitution. It does not establish the stronger universal implication
$\mathcal L(f^m)=0$ for all $m\ge1\Rightarrow f=0$.

For a homogeneous polynomial $f$ of degree $D$ in $r$ variables,
\[
 \mathcal L(f^m)=(r)_{Dm}\E_{\mathrm{Dir}(1,\ldots,1)}[f(T)^m]. \tag{3}
\]
Restriction to $\sum_iT_i=1$ and homogenization are inverse correspondences
when the degree is retained. Therefore the homogeneous Factorial Conjecture
in $r$ variables is equivalent to absence of nonzero pure-null polynomial
functions for uniform integration on the $(r-1)$-simplex. The first
two-dimensional simplex target is homogeneous FC in three variables.
The transfer itself is classical in mechanism and does not settle that target.

Relevant primary sources are
\href{https://arxiv.org/abs/1008.3962}{van den Essen--Wright--Zhao,
On the Image Conjecture} and
\href{https://arxiv.org/abs/1304.3956}{Edo--van den Essen,
The Strong Factorial Conjecture}.
The latter connects a \emph{stronger} factorial statement to Furter's
rigidity problem for gaps in coefficients of polynomial inverses.
An ordinary pointwise safety decision is not a proof of that stronger
consecutive-moment assertion.

Current status matters: Long's
\href{https://github.com/long-mathematics/factorial-conjecture-two-variables}{current FC(2) repository}
claims the full inhomogeneous two-variable theorem. Its README explicitly
marks the manuscript unpublished and unrefereed. The general flat case
uses a separate single-integral nonvanishing theorem from EIT I; the
nonflat result is stated independent of EIT I. That proof was not
independently audited here, and FC(2) should not be advertised as an
untouched target. The established Liu--Sun homogeneous binary result
also predates this assessment.

\section{A direct outward consequence for knot complexity}

The independently written quadratic-refinement argument in
\href{BREAKTHROUGH_SEARCH_TOPOLOGICAL_COUPLING_2026_11_16.tex}{the preceding topological coupling note}
has the following immediate knot-theoretic corollary.

\begin{proposition}
Suppose a knot group has a representation whose image acts transitively on
$V=\mathbb F_2^k$, preserving a common nondegenerate alternating form in
its linear parts, and every meridian maps to an affine transvection
$x\mapsto x+(B(u,x)+a)u$. Then its meridional rank, and hence its bridge
number, are at least $k+1$.
\end{proposition}
\begin{proof}
The images of a meridional generating set generate a transitive group.
Their directions must span $V$. If there are only $k$, they form a basis
and admit a quadratic refinement with $q(u_i)=1+a_i$; every generator
preserves its two level sets, a contradiction. The bridge number bounds
meridional rank from above by the standard bridge presentation.
\end{proof}

This is conditional on the existence and stated transitivity of the
representation. A linear symplectic representation fixes zero and does
not meet that hypothesis. No explicit new knot family or instance of the
meridional-rank conjecture is resolved here. Finite quotient methods
already play this role in
\href{https://doi.org/10.1007/s00208-020-02124-z}{Baader--Blair--Kjuchukova,
Coxeter groups and meridional rank of links}. The exact priority of the
affine criterion has not been established.

\section{Other transfers, with their missing steps}

\paragraph{Deformations and commutative algebra.}
The retained formal-symbol theorem states that all-marker eventual safety
over an arbitrary commutative $\mathbb Q$-algebra is equivalent to safety
after reduction, and to safety at every prime residue field; it yields an
affine multiplier-degree cutoff. This already supplies information about
infinitesimal coefficient deformations. It does not say that sharp
cutoffs, normalization ideals, or arbitrary deformation problems are
insensitive to nilpotents. The first-survival theorem also gives rational
slopes with bounded error for a specified apolar operator-power sequence,
not for all Hilbert functions or all ideal filtrations.

\paragraph{Dynamical systems and center-focus problems.}
Moment identities and Abelian periods govern first-order return
obstructions. Pakovich's
\href{https://arxiv.org/abs/1407.0150}{parametric Abel-center theorem}
demonstrates a precise moment/composition link. Actual-period quotients,
boundary extensions and exact jet selection are plausible tools for
remaining center-locus questions. Higher return obstructions involve
iterated integrals; the present Gaussian algorithm does not automatically
decide their vanishing. This is a proposed methodological transfer.

\paragraph{Dixmier and quantization.}
A genuine proof of the plane Jacobian Conjecture would imply the rank-one
Dixmier Conjecture by the published implication
$\mathrm{JC}_{2n}\Rightarrow\mathrm{DC}_n$ of
\href{https://arxiv.org/abs/math/0512171}{Belov-Kanel--Kontsevich}
(and related work of Tsuchimoto).
A partial trace theorem, a Gaussian decision algorithm, or closure of
selected HC4 graph classes does not satisfy that premise.

\paragraph{Priority.}
The most immediate exact outward application is certificate-producing
integration for Gamma and simplex laws. The closest named conjecture
family is the Factorial/Image/rigidity programme, with its stronger
universal statements retained as separate obligations. The knot bound is
another rigorous conditional transfer. Center-focus and general
quantization applications require further bridges.
\end{document}
